\section{Conclusion}
\label{sec:conclusion}

Connectivity-based autism classification on ABIDE-I is evaluated using a unified functional-connectivity pipeline. The primary ensemble achieves a participant-pooled AUROC of 0.776, with comparable discrimination retained when individual acquisition sites are held out from training. Its accuracy also reaches the range reported in previous ABIDE-I studies.

Several findings emerge from the subsequent analyses. Discrimination is weakest among autistic participants with higher cognitive ability, showing that performance varies with sample composition. Reduced somatomotor connectivity persists and becomes stronger after within-site matching for head motion. Connection-level feature selection performed before data splitting produces a large increase in reported performance, whereas region-level selection has a much smaller effect. Linear classifiers also perform better than the evaluated fixed-configuration tree ensembles on the same connectivity features and folds. In addition, classifier scores are associated with symptom measures across the pooled sample but not with symptom severity within autistic participants after correction. Finally, the apparent advantage of windowed connectivity is largely accounted for by additional covariance shrinkage, as the difference disappears when the same shrinkage is applied to the whole-scan representation. Together, these results show that ABIDE-I classification performance is strongly influenced by cohort composition and analytical choices, while also identifying a somatomotor connectivity difference that persists after motion matching.

